% =====================================================================
% SECTION 5: The arithmetic clock  +  SECTION 6: Thermodynamics & gravity
% =====================================================================

\section{The Arithmetic Clock and the Problem of Time}\label{sec:time}

\subsection{The problem of time, stated precisely}

The reconciliation problem for the two classical concepts of time is standard: in
quantum mechanics, $t$ is an external parameter and the Schr{\"o}dinger equation evolves
states with respect to it; in general relativity, elapsed time is the proper time along a
worldline, $\tau_{\gamma}=c^{-1}\int_{\gamma}\sqrt{-g_{\mu\nu}dx^{\mu}dx^{\nu}}$, a geometric
quantity that varies from trajectory to trajectory, while the canonical Hamiltonian of
gravity is a constraint, $\Hmax\Psi=0$, with no external evolution at all. The
resolution is relational: time is a correlation between subsystems, realized quantum
mechanically by the Page--Wootters mechanism \cite{pagewootters1983}, with the primon
spectrum supplying the clock. This section turns that resolution into theorems, corrects
one sign-of-life issue (the clock kernel never decorrelates, so the arrow
of time cannot live in the clock), and quantifies the claim that ``at the
Janus point time does not exist.''

The Page--Wootters construction is standard: for a total state $\Psi$ satisfying a constraint
$(H_{C}+H_{S})\Psi=0$ of clock-plus-system, the conditional state of the system when the
clock reads $\tau$ is $\psi_{S}(\tau)=\bra{\tau}\Psi\rangle$, and for an ideal clock family
this satisfies an effective Schr{\"o}dinger equation in $\tau$. All content therefore
concentrates in the clock family. The analysis below supplies the \emph{kinematical} level of
this program: the clock family, the conditional-state map, and the correlation kernel. The
constraint equation itself is not solved here---solving it requires the interaction sector of
Section~\ref{sec:dynamics} and belongs to the content of $\mathsf{C4}$---and no claim is made
that this section dissolves the problem of time; what is proved is the exact correlation
structure of the arithmetic clock, and the exact statement of what that structure forbids.

\subsection{The zeta clock}

\begin{definition}[Zeta clock family]\label{def:clock}
For a dimensionless regulator $\nu>\tfrac12$ define the normalized clock states
\begin{equation}\label{eq:clockstates}
\ket{\tau;\nu}\;=\;\zeta(2\nu)^{-1/2}\sum_{n=1}^{\infty}n^{-\nu-i\omega_{0}\tau}\ket{n},
\end{equation}
a unit-norm family in $\Hil_{\mathcal{A}}$ (unit norm because $\sum_{n}n^{-2\nu}=\zeta(2\nu)$
converges for $2\nu>1$). The regulator is a resolution parameter, not a temperature: the
weights $n^{-\nu}$ confine the clock to energies of order
$\hbar\omega_{0}/(2\nu-1)$, as computed in Proposition~\ref{prop:uncertainty}. A clock at
regulator $\nu$ has the same mean energy as a canonical primon state at $\sigma=2\nu$; the
clock family therefore exists on \emph{both} sides of the Hagedorn scale---$\nu\downarrow\tfrac12$
is the high-energy (canonical-divergence) limit, $\nu\to\infty$ the vacuum limit---and it is
the analytic object through which energies above $T_{H}$, where no Gibbs state exists, are
nevertheless resolved (Section~\ref{sec:arithmetic}).
\end{definition}

\begin{proposition}[Clock covariance and two-time kernel]\label{prop:kernel}
\begin{enumerate}[leftmargin=1.6em]
\item (Covariance.) $e^{-i\tau' H_{P}/\hbar}\ket{\tau;\nu}=\ket{\tau+\tau';\nu}$: clock
translations are generated by $H_{P}$ itself, so the clock reading $\tau$ is exactly the
phase parameter conjugate to $\log\Nhat$,
\begin{equation}
e^{-i\omega_{0}\tau\log\Nhat}\ket{n}=n^{-i\omega_{0}\tau}\ket{n}.
\end{equation}
\item (Two-time kernel.) For all $\tau,\tau'$,
\begin{equation}\label{eq:kernel}
\braket{\tau';\nu}{\tau;\nu}
\;=\;\frac{\zeta\bigl(2\nu-i\omega_{0}(\tau-\tau')\bigr)}{\zeta(2\nu)}.
\end{equation}
\item (Positivity.) The kernel \eqref{eq:kernel} is positive definite on $\Rr$: it is the
Fourier transform of the probability measure
$\mu_{\nu}=\zeta(2\nu)^{-1}\sum_{n}n^{-2\nu}\delta_{\log n}$ on the energy-per-mode
variable $u=\log n$.
\item (Product of phases.) For a superposition $\sum_{n}c_{n}\ket{n}$ evolved for time
$\tau$, the phase acquired is $n^{-i\omega_{0}\tau}$, and by unique factorization
\begin{equation}
n^{-i\omega_{0}\tau}=\prod_{p}p^{-i\omega_{0}\tau\,k_{p}}:
\end{equation}
each prime contributes an elementary phase generator, and composite evolution is the product
of prime evolutions. This is the exact sense in which ``primes provide the elementary
generators of spectral phase order.''
\end{enumerate}
\end{proposition}

\begin{proof}
(1) is immediate term by term. (2) $\braket{\tau';\nu}{\tau;\nu}
=\zeta(2\nu)^{-1}\sum_{n}n^{-2\nu}e^{i\omega_{0}(\tau-\tau')\log n}
=\zeta(2\nu)^{-1}\zeta(2\nu-i\omega_{0}(\tau-\tau'))$, since
$e^{i\omega_{0}\Delta\tau\log n}=n^{i\omega_{0}\Delta\tau}$. (3) The kernel
$K(\Delta\tau)=\int e^{i\omega_{0}\Delta\tau u}\,d\mu_{\nu}(u)$ is the characteristic
function of $\mu_{\nu}$ composed with a dilation, hence positive definite by Bochner's
theorem. (4) follows from unique factorization applied to the logarithm of the phase.
\end{proof}

Equation~\eqref{eq:kernel} is the framework's signature formula: \emph{the
distinguishability of two readings of the arithmetic clock is governed by the Riemann zeta
function}. What follows adds
the facts that make it usable.

\begin{theorem}[Almost-periodicity of the clock: the arrow cannot live here]\label{thm:bohr}
For every $\nu>\tfrac12$, the kernel $K(\Delta\tau)=\zeta(2\nu-i\omega_{0}\Delta\tau)/\zeta(2\nu)$
is a Bohr almost-periodic function of $\Delta\tau$ \cite{besicovitch1932} (the Dirichlet series
converges absolutely on the line $\sigma=2\nu>1$). Consequently:
\begin{enumerate}[leftmargin=1.6em]
\item $|K|$ does \emph{not} decay along any sequence: $\limsup_{|\Delta\tau|\to\infty}|K(\Delta\tau)|=1$.
\item For every $\varepsilon>0$ the set
$\{\Delta\tau:|K(\Delta\tau)-1|<\varepsilon\}$ is relatively dense: clock correlations
resynchronize infinitely often, in both time directions, at every resolution $\nu$.
\end{enumerate}
No limit $\nu\downarrow\tfrac12$ escapes this: at every fixed $\nu$ the kernel is two-sided
symmetric, and the singular limit only narrows the correlation peak
(Proposition~\ref{prop:uncertainty}(2)) without ever selecting a direction. The high-energy
clock is \emph{more} resolved, not \emph{directed}.
\end{theorem}

\begin{proof}
On the line $\sigma=2\nu>1$, $\zeta(s)=\sum_{n}n^{-s}$ is an absolutely convergent Dirichlet
series with $\sum_{n}n^{-\sigma}<\infty$, hence a (uniform limit of finite exponential sums)
Bohr almost-periodic function of the imaginary part \cite{besicovitch1932}. Almost-periodic
functions have relatively dense $\varepsilon$-almost-periods, and $K(0)=1$.
\end{proof}

This theorem resolves a hidden tension in the framework. The zeta
object appears twice: as the clock (``the distinguishability of two times is
governed by the zeta function'') and as the engine of the thermodynamic arrow (``the
cascade up the mountain of composite numbers''). Theorem~\ref{thm:bohr} shows
these two roles cannot be fused: an almost-periodic clock is \emph{recurrent} and carries no
arrow. In the framework the arrow is therefore relocated entirely to the
interaction sector and the boundary condition (Conjecture $\mathsf{C4}$ and $\mathsf{C7}$),
while the clock contributes the neutral, arrow-free, maximally stable phase reference. This
is a structural clarification, not a loss: it says that \emph{within the framework, time
measurement and time's arrow are functionally separated objects}, exactly as relational
philosophy of time requires.

\subsection{Energy--time duality of the arithmetic clock}

\begin{proposition}[Resolution duality]\label{prop:uncertainty}
For the clock family $\ket{\tau;\nu}$:
\begin{enumerate}[leftmargin=1.6em]
\item The mean energy is
$\langle H_{P}\rangle_{\nu}=\hbar\omega_{0}\bigl(-\zeta'/\zeta\bigr)(2\nu)
=\dfrac{\hbar\omega_{0}}{2\nu-1}+O(\hbar\omega_{0})$ as $\nu\downarrow\tfrac12$.
\item The correlation peak is asymptotically Lorentzian: as $\nu\downarrow\tfrac12$,
uniformly for bounded $\omega_{0}\Delta\tau/(2\nu-1)$,
\begin{equation}
|K(\Delta\tau)|^{2}
=\frac{(2\nu-1)^{2}}{(2\nu-1)^{2}+\omega_{0}^{2}\Delta\tau^{2}}\,\bigl(1+o(1)\bigr),
\end{equation}
with half-width $\Delta\tau_{1/2}\asymp(2\nu-1)/\omega_{0}$.
\item Consequently $\langle H_{P}\rangle_{\nu}\,\Delta\tau_{1/2}\to\hbar$ as
$\nu\downarrow\tfrac12$: the zeta clock asymptotically saturates the
Mandelstam--Tamm energy--time limit. Time resolution is bought one-for-one with clock energy,
and the clock never becomes ``cheap.''
\item Conversely, as $\nu\to\infty$ the state concentrates on $\ket{1}$:
$\ket{\tau;\nu}\to\ket{1}$, energy $\to0$, and $K\to1$ for all $\Delta\tau$: the clock
resolves nothing. At the Janus limit the clock degenerates---the quantified form of the
claim that ``at $n=1$ time does not exist.''
\end{enumerate}
\end{proposition}

\begin{proof}
(1) is Theorem~\ref{thm:hagedorn} at $\sigma=2\nu\downarrow1$. Positivity of $\langle H_P\rangle$ requires $\nu>1/2$, consistent with Definition~\ref{def:clock}. (2) Near the pole,
$\zeta(\sigma-it)=(\sigma-1-it)^{-1}+O(1)$, so
$|K|^{2}=|\sigma-1|^{2}/((\sigma-1)^{2}+t^{2})+o(1)$ with $t=\omega_{0}\Delta\tau$ and
$\sigma-1=2\nu-1$. (3) Multiply (1) and (2). (4) $n^{-2\nu}$ concentrates on $n=1$;
$K(t)=\zeta(2\nu-it)/\zeta(2\nu)\to1$ since $\zeta(2\nu-it)\to1$ and
$\zeta(2\nu)\to1$.
\end{proof}

\subsection{Reconciliation of the three times}

With the clock machinery in place, the tripartite reconciliation can be stated as
a compact conditional statement. \emph{Quantum-mechanical time} is the reading of a
relational clock in the classical-clock limit: it is recovered when the clock subsystem is
large, weakly coupled, and quasi-classical, in which case $\tau_{\mathrm{clock}}\approx t$
for all practical observers. \emph{Relativistic time} is the proper-time functional along
worldlines: it is the geometric content of the clock's embedding, determining which phase
increments accumulate along which trajectory. \emph{Quantum-gravitational time} is the
constraint-generated phase correlation: with the primon clock, $\tau$ is conjugate to
$\log\Nhat$, and evolution is the multiplicative phase flow
$n^{-i\omega_{0}\tau}$. The three are not rivals but levels: the block-level correlation
structure ($\Hmax\Psi=0$), its geometric projection (proper time), and its semiclassical
limit (external $t$). All incompatibilities of principle dissolve at the kinematical level;
what remains open
is whether the \emph{specific} identification of the clock with the
primon spectrum is realized (Conjecture $\mathsf{C1}$), and how the geometric sector
emerges from the arithmetic one (Conjectures $\mathsf{C3}$, $\mathsf{C9}$).

The tripartite reconciliation is compactly stated:
\begin{center}
\emph{Relativity provides the geometry; quantum mechanics provides the spectrum; the primes
provide the arithmetic order of the spectrum; and time is the phase relation among
prime-ordered states measured along geometric worldlines.}
\end{center}
The measurement clause is Proposition~\ref{prop:kernel}(1) (clock covariance); the ordering
clause is unique factorization, i.e. Proposition~\ref{prop:kernel}(4); the arrow clause is
conspicuously absent, by Theorem~\ref{thm:bohr}, and lives in the interaction sector.

\section{Thermodynamics, Gravity, and the Holographic Bridge}\label{sec:thermo}

\subsection{The exact counting law}

A repeatedly used statement is ``the number of states below energy $E$ grows like
$e^{E/\hbar\omega_{0}}$.'' The exact statement is a one-line identity, and its precision is
what makes the holographic bridge in Section~\ref{sec:holography} clean.

\begin{theorem}[Log-count identity]\label{thm:logcount}
Let $N(E)=\#\{n\in\Nn_{\geq1}:\hbar\omega_{0}\log n\leq E\}=\lfloor e^{E/\hbar\omega_{0}}\rfloor$,
and write $x=E/\hbar\omega_{0}$ and $r(E)=\log N(E)-x$. Then:
\begin{enumerate}[leftmargin=1.6em]
\item (Global envelope.) For all $E\geq0$: $-\log2\leq r(E)\leq0$.
\item (Sharp two-sided bound, $E\geq\hbar\omega_{0}\log2$.)
\begin{equation}\label{eq:sharpdefect}
\log\bigl(1-e^{-x}\bigr)\;\leq\; r(E)\;\leq\; 0,
\end{equation}
with both bounds sharp: $r\uparrow0$ as $E$ crosses a level from above, and
$r\downarrow\log(1-e^{-x})$ as $E$ approaches a level from below.
\item (Exponential decay.) $|r(E)|\leq-\log(1-e^{-x})\leq 2e^{-x}$: the defect converges to
zero \emph{exponentially}. The counting function itself remains a staircase forever: $N$
jumps by one at every level $E=\hbar\omega_{0}\log m$, and the jump of $\log N$ there is
$\log\bigl(1+\tfrac1m\bigr)\sim\tfrac1m$---the granularity persists at every scale while its
amplitude decays like $e^{-x}$.
\item (Smoothing.) For every probability kernel $w$ supported in $[-\Delta,\Delta]$ with
$\Delta$ fixed, the smoothed count $\widetilde N=N\ast w$ satisfies, uniformly in
$E\geq\hbar\omega_{0}(\log2+\Delta)$,
\begin{equation}\label{eq:smoothing}
\log\widetilde N(E)\;=\;x\;+\;\log\!\int e^{t}w(t)\,dt\;+\;O(e^{-x}):
\end{equation}
the growth rate of every fixed-window smoothing of the entropy is exactly $1/\hbar\omega_{0}$
up to exponentially small corrections, at every energy scale.
\end{enumerate}
In particular the microcanonical entropy $S(E)=k_{B}\log N(E)$ satisfies
\begin{equation}
S(E)\;=\;\frac{E}{T_{H}}\;+\;\delta(E),
\qquad
|\delta(E)|\;\leq\;k_{B}\log\frac{1}{1-e^{-E/\hbar\omega_{0}}}\;\leq\;2k_{B}e^{-E/\hbar\omega_{0}}
\quad(E\geq\hbar\omega_{0}\log2),
\end{equation}
with $T_{H}=\hbar\omega_{0}/k_{B}$. (The unsmoothed entropy is a staircase; its pointwise
microcanonical temperature is undefined between levels. The content of the Hagedorn reading
is item~(iv): the growth rate is energy-independent, so energy input converts into state
multiplicity at a fixed exchange rate.)
\end{theorem}

\begin{proof}
$N(E)=\lfloor e^{x}\rfloor$ because $\hbar\omega_{0}\log n\leq E\iff n\leq e^{x}$. Write
$y=e^{x}$ and $\{y\}=y-\lfloor y\rfloor\in[0,1)$, so $r=\log\bigl(1-\{y\}/y\bigr)\in[\log(1-1/y),0]$
since $\lfloor y\rfloor\geq y-1$ and $\{y\}<1$. The lower bound is approached as $\{y\}\uparrow1$
($E$ approaching a level from below), the upper bound attained as $\{y\}\downarrow0$. Decay:
$-\log(1-u)\leq2u$ for $u\leq\tfrac12$, and $\{y\}/y<1/y=e^{-x}\leq\tfrac12$ for $x\geq\log2$.
Jumps: on $y\in[m,m+1)$, $\log N=\log m$ and the jump at $y=m+1$ is
$\log\bigl((m+1)/m\bigr)\sim1/m$. Smoothing: $N(E+t)=e^{x+t}-\{e^{x+t}\}$, so
$\widetilde N(E)=e^{x}\int e^{t}w(t)\,dt-R(E)$ with $0\leq R(E)=\int w(t)\{e^{x+t}\}\,dt\leq1$;
hence $\log\widetilde N=x+\log\int e^{t}w(t)dt+\log\bigl(1-R(E)e^{-x}/\!\int e^{t}w\bigr)$ and
the last term is $O(e^{-x})$, uniformly in $E$.
\end{proof}

Two consequences of the corrected defect estimate are load-bearing later, and both follow
from item~(iii). First, at matched entropy scale the integer-counting fluctuation is
$O(k_{B}e^{-E/\hbar\omega_{0}})$: \emph{far smaller} than the logarithmic corrections of
black-hole entropy, so those corrections cannot originate from the floor sector at
all---Section~\ref{sec:holography} turns this from a caveat into a falsifier (the counting
contributes no $\log A$ term; any such term must be carried by the dictionary). Second, the
fluctuation sector of the gravitational counting function is carried not by the floor but by
the \emph{prime}-counting oscillations: the explicit formula
$\psi(x)=x-\sum_{\rho}x^{\rho}/\rho+\cdots$ has fluctuations bounded by $O(\sqrt x\log^{2}x)$
under the Riemann hypothesis---the content of von Koch's equivalence
\cite{vonkoch1901}: the Riemann hypothesis is \emph{equivalent} to the statement that
arithmetic fluctuations about the smooth classical term are as small as
square-root-small. Read in the framework's dictionary, RH becomes: \emph{quantum
fluctuations of the arithmetic sector about the classical continuum are optimally
suppressed}. That is a rigorous reformulation of ``RH = spectral
stability'', and it is the strongest form of that claim that mathematics currently
licenses.

\subsection{The holographic bridge}\label{sec:holography}

The derivation of the Bekenstein--Hawking entropy from prime state counting
is a natural target. This paper isolates exactly one
nontrivial input and shows the rest is Theorem~\ref{thm:logcount}.

\begin{conjecture}[Arithmetic holographic bound]\label{conj:holobound}
The maximal primon excitation energy accessible in, or associated with, a spacetime region
whose boundary has area $A$ is
\begin{equation}\label{eq:holobound}
E_{\max}(A)\;=\;\frac{\hbar\omega_{0}}{4\ell_{P}^{2}}\,A\;+\;\hbar\omega_{0}\,s(A),
\qquad
s(A)=O\!\Bigl(\log\frac{A}{\ell_{P}^{2}}\Bigr),
\end{equation}
independently of the detailed matter content, where
$\ell_{P}=\sqrt{G\hbar/c^{3}}$ is the Planck length. The term $\hbar\omega_{0}s(A)$ is the
\emph{dictionary error term}: it is the only slot through which logarithmic corrections can
enter, since the counting theorem contributes none (Theorem~\ref{thm:logcount}(iii)).
\end{conjecture}

\begin{theorem}[Conditional Bekenstein--Hawking]\label{thm:bh}
Assume Conjecture~\ref{conj:holobound}. Then the arithmetic entropy of the region is
\begin{equation}
S(A)\;=\;k_{B}\log N\bigl(E_{\max}(A)\bigr)
\;=\;\frac{k_{B}A}{4\ell_{P}^{2}}\;+\;k_{B}\,s(A)\;+\;\delta\bigl(E_{\max}(A)\bigr),
\qquad
|\delta|\leq2k_{B}e^{-E_{\max}/\hbar\omega_{0}},
\end{equation}
which is the Bekenstein--Hawking entropy \cite{bekenstein1973,hawking1975} with the area
coefficient exact and \emph{every} logarithmic correction carried by the dictionary term
$s(A)$ and none by the counting. The pure model ($s=O(1)$) predicts no
$-\alpha\log(A/\ell_{P}^{2})$ term at all; a computed correction of that class (several
ensembles give $-\tfrac32\log A$) therefore falsifies the pure form and measures $s$.
Moreover, for a Schwarzschild black hole of
mass $M$, $A=16\pi G^{2}M^{2}/c^{4}$ and $A/4\ell_{P}^{2}=4\pi GM^{2}/(\hbar c)$; the choice
$E_{\max}=Mc^{2}$ and $\hbar\omega_{0}=2\,k_{B}T_{H}^{\mathrm{BH}}$ normalizes
\eqref{eq:holobound} \emph{exactly} at fixed horizon, with $T_{H}^{\mathrm{BH}}=\hbar c^{3}/(8\pi GMk_{B})$
the Hawking temperature.
\end{theorem}

\begin{proof}
Substitute \eqref{eq:holobound} into Theorem~\ref{thm:logcount}. For the Schwarzschild case,
$E_{\max}(A)/\hbar\omega_{0}=A/4\ell_{P}^{2}$ by \eqref{eq:holobound}; with
$E_{\max}=Mc^{2}$ this forces $\hbar\omega_{0}=4Mc^{2}\ell_{P}^{2}/A
=4Mc^{2}\cdot(G\hbar/c^{3})/(16\pi G^{2}M^{2}/c^{4})
=\hbar c^{3}/(4\pi GM)=2k_{B}T_{H}^{\mathrm{BH}}$.
\end{proof}

Two remarks delimit the bridge; both are registered under $\mathsf{C3}$.

\begin{remark}[The factor-of-two ambiguity]
The Hawking identification above is exact only per-horizon: since $T_{H}^{\mathrm{BH}}\propto
M^{-1}$, a \emph{universal} $\omega_{0}$ cannot equal $2T_{H}^{\mathrm{BH}}$ for all masses.
Either $\omega_{0}$ is universal and the identification holds to within
$O(1)$ horizon-dependent factors, or the arithmetic scale is set per horizon. Which
alternative survives is a genuine open question that the conjecture form makes explicit;
a presentation that silently treats $\omega_{0}$ as universal and the
identification as exact would hide it. In both alternatives the \emph{area coefficient} is
exact, because the counting defect is exponentially small (Theorem~\ref{thm:logcount}(iii));
all residual freedom lives in the dictionary term $s(A)$.
\end{remark}

\begin{remark}[What the bridge does and does not derive]
The bridge converts an area law into an \emph{energy} cutoff statement plus an exact
counting theorem; it does not derive the cutoff. This is precisely the logical position of
the holographic principle in mainstream theory, translated into arithmetic terms. The
conversion loss is provably zero in the area coefficient (Theorem
\ref{thm:logcount} has no hidden degeneracy factors, since the primon spectrum is simple,
Theorem~\ref{thm:spectrum}), and every logarithmic correction is imported through the
dictionary term $s(A)$: any slack in Conjecture~\ref{conj:holobound} is slack in the
physics, not in the counting.
\end{remark}

\subsection{Singularities and the Hagedorn wall}

The singularity claims swing between two poles: one places the primon gas
\emph{at} the singularity (the reported conformal system near the black-hole center), later
reading declares that singularities \emph{do not exist}. The reconciliation is a standard
move of effective field theory, and it is worth stating once, cleanly, because it removes
the apparent contradiction:

\begin{remark}[Resolution of the singularity tension]\label{rem:singularity}
``Singularity'' denotes the boundary of validity of the continuum effective description, not
a locus in the fundamental theory. In the reconstructed framework: the geometric (Einstein)
description fails where its implicit continuum counting diverges; the arithmetic description
remains finite at all energies (Theorem~\ref{thm:spectrum} has no finite accumulation
point, and all quantities are finite below $T_{H}$). The pole of $\zeta$ at $s=1$ is the
exact analytic signature of the breakdown: it is a \emph{critical} (Hagedorn) divergence,
not a curvature divergence. What the Big Bang and black-hole interiors ``are'' in this
framework is the regime where the geometric effective theory must be replaced by the
arithmetic one; the primon gas is the UV-complete description of \emph{that regime}, which is
compatible with it also being located, in the geometric coordinates of the effective theory,
``near the singularity.'' Both readings are thereby retained, each in its
proper tense. The associated physical conjecture ($\mathsf{C6}$) is that the Hagedorn wall
$T_{H}$ replaces the singular point: approach to the singularity is approach to a critical
temperature at which energy converts into primon creation rather than curvature growth.
\end{remark}

\subsection{Emergent gravity and the hierarchy problem}

The resolution of the hierarchy problem---gravity as an entropic, $1/N$-suppressed
collective effect \cite{jacobson1995,verlinde2011,thooft1974}---is the most speculative of
its physics claims, and it is demoted in full. But the demoted form can be made
unusually concrete, because the arithmetic sector supplies an explicit $N$.

\begin{conjecture}[Arithmetic hierarchy]\label{conj:hierarchy}
The Newton constant is set by the number of arithmetic degrees of freedom below the
effective cutoff $E_{\mathrm{cut}}$:
\begin{equation}
\frac{1}{G}\;\asymp\;\frac{\hbar}{c^{3}}\,\omega_{0}\,
N_{\mathrm{eff}}\bigl(E_{\mathrm{cut}}\bigr),
\qquad
N_{\mathrm{eff}}(E)=\lfloor e^{E/\hbar\omega_{0}}\rfloor,
\end{equation}
up to a dimensionless $O(1)$ factor, so that
\begin{equation}\label{eq:hierlog}
\log\frac{M_{\mathrm{Pl}}}{M_{\mathrm{EW}}}\;\asymp\;
\frac{E_{\mathrm{cut}}}{2\hbar\omega_{0}}
\qquad\Longrightarrow\qquad
E_{\mathrm{cut}}\;\asymp\;2\log\!\Bigl(\frac{M_{\mathrm{Pl}}}{M_{\mathrm{EW}}}\Bigr)\,\hbar\omega_{0}
\;\approx\;74\,\hbar\omega_{0}
\end{equation}
(using $M_{\mathrm{Pl}}/M_{\mathrm{EW}}\approx10^{16}$).
\end{conjecture}

The honest content of Conjecture~\ref{conj:hierarchy} is a \emph{relocation} of the hierarchy
problem, stated with numbers: the question ``why is gravity so weak?'' becomes ``why does the
effective arithmetic universe of the electroweak regime terminate at $n\sim e^{74}\approx
10^{32}$ integers?''---which is a question about the cutoff structure of the arithmetic
sector, not a fine-tuning of a Lagrangian parameter. The framework does not answer it; it
converts it into a well-posed question with a single dimensionless number to explain,
register entry $\mathsf{C11}$. This is strictly more than a bare
$1/N$-suppression narrative with no $N$, and strictly less than a derivation of $G$.

\subsection{The classical limit and the pole at unity}

The classical-limit condition reads: as $s\to1^{+}$, the coarse-grained
arithmetic partition function should reproduce classical geometry, the pole
$\zeta(s)\sim1/(s-1)$ representing the divergent continuum limit. This paper
restates this as a theorem-flanked analogy: the pole is exactly the Hagedorn divergence
(Theorem~\ref{thm:hagedorn}); the smooth term $x$ of $\psi(x)$ is the classical contribution
and the zeros contribute the fluctuations (von Koch's equivalence
\cite{vonkoch1901}); and ``renormalization removing the pole'' means passing to the
finite part of the Laurent expansion, whose constant term is Euler's constant $\gamma$---the
arithmetic framework's first finite counterterm. Whether the geometric sector (Einstein
dynamics, $G_{\mu\nu}+\Lambda g_{\mu\nu}=8\pi GT_{\mu\nu}$) can actually be \emph{derived} as
this coarse-graining is the deepest part of Conjecture $\mathsf{C1}$ and is not claimed.
